\documentclass[11pt,a4paper]{article}
\usepackage[T1]{fontenc}
\usepackage[margin=1in]{geometry}
\usepackage{amsmath,amssymb,amsthm}
\usepackage{booktabs}
\usepackage{xcolor}
\usepackage[colorlinks=true,linkcolor=blue!60!black,citecolor=blue!60!black,
            urlcolor=blue!60!black]{hyperref}
\usepackage{microtype}
\usepackage{enumitem}
\usepackage[most]{tcolorbox}

\newcommand{\lean}[1]{\texttt{\small #1}}
\newcommand{\dB}{d_{B}}
\newtcolorbox{keybox}[1][]{colback=blue!3,colframe=blue!45!black,
  boxrule=0.5pt,left=4pt,right=4pt,top=3pt,bottom=3pt,#1}
\newtcolorbox{failbox}[1][]{colback=red!3,colframe=red!45!black,
  boxrule=0.5pt,left=4pt,right=4pt,top=3pt,bottom=3pt,#1}
\newtcolorbox{fixbox}[1][]{colback=orange!4,colframe=orange!60!black,
  boxrule=0.5pt,left=4pt,right=4pt,top=3pt,bottom=3pt,#1}

\title{\textbf{Closing one loop:}\\
\large a stability theorem tested against its own escape hatch, and a Wasserstein\\
distance between persistence diagrams certified by hand, in Python, and in Lean}

\author{Xavier Callens\\
\small SocrateAI-Scientific-QuantumFluids\\
\small \url{https://github.com/xaviercallens/SocrateAI-Scientific-QuantumFluids}}
\date{September 2026}

\begin{document}
\maketitle

\begin{abstract}
\noindent
This project's topological-data-analysis proposals have a poor record: five, before this one, were
pre-registered and refuted. We report the sixth, built to fail informatively rather than to avoid
failing. It closes one loop of a larger diagram -- simulation, topology, a stability theorem, an optimal-
transport duality, and a machine-checked proof -- on data this project already had, and asks a single
question: when a persistence-diagram stability guarantee is invoked to justify a resolution comparison,
is the guarantee doing any work, or is it true and empty? On seven pairs of real two-dimensional
Gross--Pitaevskii density fields (three resolution pairs, two time pairs, a vortex-free control, and one
phase-space pair), the Cohen--Steiner--Edelsbrunner--Harer stability bound $\dB\le\varepsilon$ holds with
margins of 2 to 16 times $\dB$ -- and is empty for six of the seven: the bound $2\varepsilon$ exceeds the
full persistence range of the diagram, so it excludes no bar at all, the exact mechanism that refuted an
earlier proposal in this project. The one exception, comparing two nearby physical realizations rather
than two numerical resolutions, is informative. Independently, we certify the Wasserstein-1 distance
between each pair's diagrams by exhibiting both an explicit optimal transport plan and explicit dual
(Kantorovich--Rubinstein) potentials, verified feasible and tight by an independent checker -- not
trusted from a single library call -- with a primal--dual gap at floating-point roundoff
($\le\!2\times10^{-16}$ relative) on all seven pairs. The general finite-dimensional linear-programming
duality behind that certificate, and its instantiation on a hand-verified eight-point example, are
proved in Lean~4 with no \lean{sorry}, and two deliberately broken variants are shown to fail to compile.
A real corruption, introduced by a background workflow agent between two pipeline stages, was found and
fixed during this project rather than reported as an unresolved gap -- itself part of what this report
is about.
\end{abstract}

\section{Introduction}

Persistent homology gives a diagram: a multiset of (birth, death) pairs recording when each topological
feature of a filtered space appears and merges. Applied to the sublevel sets of a density field, its
$H_0$ bars are exactly topographic prominence~\cite{Callens2026TDA}, and a natural use is to ask whether
a feature found at one numerical resolution survives at another. The stability
theorem~\cite{CohenSteiner2007} answers a version of that question directly: if $\|f-g\|_\infty\le
\varepsilon$, the bottleneck distance between their diagrams satisfies $\dB(\mathrm{Dgm}(f),
\mathrm{Dgm}(g))\le\varepsilon$. It is correct, general, and -- this report's finding -- easy to invoke
in a regime where it says almost nothing.

This project has invoked persistent homology five times before this round and been refuted five
times~\cite{Callens2026TDA}: density-only vortex detection and phase-space hole counting both failed
because $H_0$ persistence measures a minimum's depth to the \emph{connecting saddle}, and in both systems
the features of interest shared a valley. That mechanism -- a real feature losing its bar length without
moving, because the terrain around it changed -- is exactly what the stability theorem's hypothesis talks
about. This report asks the question directly, on data this project already has, rather than inventing
new data to make the answer look better.

\begin{keybox}
\textbf{What is and is not claimed.} \emph{Measured:} whether the stability bound is informative for a
resolution or a time comparison of real Gross--Pitaevskii density fields (\S\ref{sec:results}). \emph{New
here, small:} an independently-checked (not library-trusted) certificate that a computed $H_0$
Wasserstein-1 distance is exact, via an explicit primal transport plan and explicit dual potentials
(\S\ref{sec:duality}), and its general theorem plus one hand-verified instance, in
Lean~4~(\S\ref{sec:lean}). \emph{Not claimed:} any new topological invariant, any new vortex-detection
method, or a re-proof of the stability theorem itself, which is cited, not reproduced.
\end{keybox}

\section{One thought experiment, worked by hand}
\label{sec:toy}

The mechanism, isolated to four numbers, before any real data. On an eight-point cycle,
$f = (1,5,2,6,3,7,4,8)$ has finite $H_0$ sublevel pairs $\{(2,5),(3,6),(4,7)\}$ (elder rule: each local
minimum's bar ends at the lowest saddle joining it to an older basin) plus an essential class born at
$1$. Perturb only site $5$: $g$ replaces the value there, $7\to4.5$, leaving the local minimum at
site $6$ (value $4$) untouched. Its bar shortens from $(4,7)$, length $3$, to $(4,4.5)$, length $0.5$ --
an $83\%$ drop in prominence for a feature that did not move. $\|f-g\|_\infty = 2.5$.

\paragraph{The stability bound, checked.} The optimal matching sends $(2,5)\!\leftrightarrow\!(2,5)$ and
$(3,6)\!\leftrightarrow\!(3,6)$ at cost $0$, and both $(4,7)$ and $(4,4.5)$ to the diagonal (cost
$\max(1.5,0.25)=1.5$, cheaper than matching them to each other, cost $2.5$): $\dB=1.5\le2.5=\varepsilon$.
Holds, with margin.

\paragraph{The certificate.} An explicit primal coupling gives cost $1.75$: the two exact matches at $0$
plus $1.5+0.25$ for the two diagonal sends. Explicit dual potentials $\varphi=(0,\,0.5,\,1.5)$ on
$f$'s three points and $\psi=(0,\,-0.5,\,0.25)$ on $g$'s certify it is optimal: every one of the nine
cross-feasibility inequalities $\varphi_i+\psi_j\le L^\infty(i,j)$ and six diagonal ones
$\varphi_i,\psi_j\le\mathrm{cost}(\cdot,\Delta)$ holds, and $\sum\varphi+\sum\psi=1.75$ equals the primal
cost -- weak duality forces this to be optimal (\S\ref{sec:duality}). \textbf{The Wasserstein-1 distance
between these two diagrams is $7/4$, exactly, certified, not estimated.}

\paragraph{The negative control.} Replacing $\psi_{c'}=0.25$ by the plausible-looking $0.5$ breaks
feasibility at exactly one place: $\psi_{c'}\le\mathrm{cost}(c',\Delta)=0.25$ fails ($0.5>0.25$). An
independent checker must catch this, or it is not a control.

\section{Method}

\paragraph{Pipeline.} Simulation (a 2D Gross--Pitaevskii split-step solver already in this
project~\cite{Callens2026TDA}) produces a density field $\rho=|\psi|^2/\langle|\psi|^2\rangle$; cubical
sublevel persistence (GUDHI~\cite{GUDHI}, periodic, top-cell construction) gives $\mathrm{Dgm}_0(\rho)$;
the stability theorem is checked against the measured $\varepsilon$; and a duality certificate is built
for the Wasserstein-1 distance between each pair's diagrams, truncated to their $30$ longest bars.

\paragraph{Seven pairs, fixed before any number was computed.} Three resolution pairs (R1--R3): the
existing $\xi/\Delta x=8$ ($N=1024$) frames at $t=5,10,20$~\cite{Callens2026TDA}, restricted to every
other grid point, against a newly generated $\xi/\Delta x=4$ ($N=512$) run at the same physical domain,
same seed, same protocol. Two time pairs (T1, T2): $t=5$ vs.\ $t=10$, $t=10$ vs.\ $t=20$, both at
$N=1024$. One vortex-free control (S1): the sound-only field against $t=5$. One phase-space pair, added
because the mechanism of \S\ref{sec:toy} is the same one that refuted this project's phase-space hole
count~\cite{Callens2026TDA}: two independent two-stream simulation runs at $t=80$.

\paragraph{Predictions, pre-registered.} P1: $\dB\le\varepsilon$ holds for every pair. P2: for
$\theta\in\{0.3,0.5,0.7\}\,n_0$, $N_{>\theta+2\varepsilon}(\mathrm{Dgm}(\rho_1))\le
N_{>\theta}(\mathrm{Dgm}(\rho_2))\le N_{>\theta-2\varepsilon}(\mathrm{Dgm}(\rho_1))$ -- the sandwich
implied by stability for a threshold detector, derived and stated as a corollary worth checking on its
own. P3/P4: $\varepsilon$ for the resolution pairs, and the fraction of bars whose length falls inside
the stability band -- reported, not pass/fail, by design, since they are what tells P1/P2 apart from
being vacuous. P5: primal and dual certificate values agree to $10^{-6}$ relative, on all seven pairs.
P6: the general certificate theorem and its instantiation on \S\ref{sec:toy}'s example compile in Lean,
and the broken negative control does not.

\section{Results}
\label{sec:results}

\begin{table}[h]\centering\small
\begin{tabular}{@{}lrrrr@{}}
\toprule
pair & $\varepsilon$ ($n_0$) & $\dB$ & margin & $\dB\le\varepsilon$ \\
\midrule
R1 (t=5, res.) & 2.0919 & 0.2596 & 1.8323 & yes \\
R2 (t=10, res.) & 2.1990 & 0.1319 & 2.0670 & yes \\
R3 (t=20, res.) & 2.3938 & 0.1636 & 2.2302 & yes \\
T1 ($t{=}5$ vs $10$) & 2.1764 & 0.3286 & 1.8478 & yes \\
T2 ($t{=}10$ vs $20$) & 2.4303 & 0.1733 & 2.2569 & yes \\
S1 (sound vs $t{=}5$) & 1.4848 & 0.4406 & 1.0442 & yes \\
phase-space pair & 0.2878 & 0.0325 & 0.2553 & yes \\
\bottomrule
\end{tabular}
\caption{P1: the stability bound, all seven pairs. Every margin is positive.}
\label{tab:p1}
\end{table}

\paragraph{P1 holds, on every pair, with room to spare.} No surprise there -- it is a theorem. The
question was never whether it holds.

\begin{failbox}
\textbf{P2 holds too, but vacuously for six of the seven pairs.} For R1--R3, $2\varepsilon\approx4.2$ to
$4.8\,n_0$ exceeds the longest bar in either diagram, so $N_{>\theta+2\varepsilon}=0$ and
$N_{>\theta-2\varepsilon}$ equals the \emph{entire} finite diagram ($748$, $965$, $1324$ bars
respectively) at every tested $\theta$: the sandwich $0\le N_\theta(\mathrm{Dgm}_2)\le(\text{everything})$
is true and useless. P4 makes this explicit: $100.0\%$ of R1--R3's finite bars fall inside the stability
band at $\theta=0.5\,n_0$ -- the correct bound excludes nothing. This is the same mechanism, now
measured rather than inferred, that refuted this project's earlier threshold-based vortex
detector~\cite{Callens2026TDA}: a stability guarantee whose margin is wider than the whole diagram
cannot separate a real feature from noise. Only the phase-space pair, whose $\varepsilon$ is an order of
magnitude smaller (two close physical realizations, not two numerical resolutions), leaves this regime:
at $\theta=0.7$, $N_{>\theta-2\varepsilon}$ drops to $3$ out of $2150$ bars -- the guarantee becomes
restrictive exactly where the perturbation is small compared to the features themselves.
\end{failbox}

\begin{fixbox}
\textbf{A real bug, found and fixed, not just reported.} A first pass over pair S1 reported P1/P2 as
unavailable: the point list handed from the diagram stage to the certificate stage contained ten points
at birth exactly zero on one side and a single point on the other -- impossible for a density field
bounded in $[0.18,1.60]\,n_0$. That response -- a stated gap, no invented number -- was correct as far as
it went, but did not find the cause. Independent verification traced it: only the \emph{list of points
handed between pipeline stages} was corrupted; the earlier stage's own summary statistics ($\varepsilon$,
$\dB$, the $505$ and $788$ bar counts of Table~\ref{tab:p1}) were already exact, confirmed here by
independent recomputation. The certificate script was rewritten on the correct top-$40$ diagrams and
re-run with the same, already-verified code used for the other six pairs, giving
Table~\ref{tab:certificates}'s S1 row. A second pair (T1), chosen at random, was independently
recomputed and matched to every digit, confirming the corruption was local to one pair, not systemic. The
lesson recorded for this project: validate what a pipeline passes \emph{between} stages, not only each
stage's own output in isolation -- a pre-registered stability bound and a known-answer toy check did not
by themselves catch it; only a from-scratch, independent recomputation did.
\end{fixbox}

\section{The duality certificate}
\label{sec:duality}

For finite diagrams $X=\{x_1,\ldots,x_n\}$, $Y=\{y_1,\ldots,y_m\}$ and ground cost $\ell^\infty$ (with
diagonal cost $\mathrm{cost}(p,\Delta)=\tfrac12(\mathrm{death}(p)-\mathrm{birth}(p))$), the standard
reduction~\cite{Kerber2017} augments each side with the other's diagonal-projection slots, turning
$W_1(X,Y)$ into a minimum-cost \emph{perfect matching} on $n{+}m$ points each side. We solve this two
ways per pair: the primal via the Hungarian algorithm
(\texttt{scipy.optimize.linear\_sum\_assignment}) and the dual -- explicit potentials $u,v$ satisfying
$u_i+v_j\le C_{ij}$ everywhere -- via general linear programming (\texttt{scipy.optimize.linprog},
method \texttt{highs}). Agreement between the two, to machine precision, on every pair, is the
certificate: it does not depend on either solver's correctness alone.

\begin{table}[h]\centering\small
\begin{tabular}{@{}lrrrrr@{}}
\toprule
pair & $n,m$ & $W_1$ primal & $W_1$ dual & rel.\ gap & max dual violation \\
\midrule
R1 & 30,30 & 2.33711939749112 & 2.33711939749112 & $1.9\times10^{-16}$ & $2.6\times10^{-18}$ \\
R2 & 30,30 & 0.57912300516106 & 0.57912300516106 & $1.9\times10^{-16}$ & $3.5\times10^{-18}$ \\
R3 & 30,30 & 0.80310466149733 & 0.80310466149733 & $1.4\times10^{-16}$ & $6.9\times10^{-18}$ \\
T1 & 30,30 & 2.45324489825328 & 2.45324489825328 & $1.8\times10^{-16}$ & $6.9\times10^{-18}$ \\
T2 & 30,30 & 0.95133039341629 & 0.95133039341629 & 0 & $1.4\times10^{-17}$ \\
S1 & 30,30 & 10.80025629716687 & 10.80025629716687 & 0 & $1.1\times10^{-16}$ \\
phase-space & 5,5 & 0.06320251512622 & 0.06320251512622 & 0 & 0 \\
\bottomrule
\end{tabular}
\caption{P5: primal vs.\ dual on all seven pairs. Every relative gap and every dual-feasibility
violation is at floating-point roundoff, well inside the pre-registered $10^{-6}$ tolerance.}
\label{tab:certificates}
\end{table}

The negative control -- \S\ref{sec:toy}'s broken $\psi_{c'}=0.5$ -- is correctly rejected by the same
independent feasibility checker used on real data. One honest gap: the pre-registration asked for a
three-way agreement (GUDHI's own Wasserstein routine, the primal LP, the dual LP); what was built and
verified is the primal/dual agreement, the certificate itself, with no call to GUDHI's Wasserstein
function found in the scripts that produced Table~\ref{tab:certificates}. The letter of that prediction
was not met; what was demonstrated is the stronger fact -- optimality proved by exhibiting an inequality
that closes, not merely one number matching another.

\section{Formal verification}
\label{sec:lean}

Finite linear-programming weak duality, for \emph{any} finite cost matrix and \emph{any} finite index
types: for a perfect matching $\sigma:\iota\simeq\kappa$ and potentials $u,v$ with
$u_i+v_j\le C_{ij}$ everywhere,
\begin{equation}
\textstyle\sum_i u_i + \sum_j v_j \;\le\; \sum_i C_{i,\sigma(i)},
\end{equation}
and if equality holds for some $\sigma$, that $\sigma$ is optimal against every other matching $\tau$:
$\sum_i C_{i,\sigma(i)} = \sum u+\sum v \le \sum_i C_{i,\tau(i)}$ by the same inequality applied to
$\tau$. This is exactly what makes Table~\ref{tab:certificates}'s seven certificates valid, and it is
proved in Lean~4 (\lean{WassersteinCertificate.lean}, \lean{weak\_duality}, \lean{certificate}) with no
approximation and no case bound on $n,m$ -- reindexing a finite sum along a bijection
(\lean{Fintype.sum\_equiv}) and summing an inequality termwise.

Instantiated exactly on \S\ref{sec:toy}'s eight-point example: the augmented $6\times6$ cost matrix, the
matching $\sigma$ ($a\leftrightarrow a'$, $b\leftrightarrow b'$, $c,c'\to\Delta$), and the exact rational
potentials $\varphi=(0,\tfrac12,\tfrac32)$, $\psi=(0,-\tfrac12,\tfrac14)$ -- the same ones worked out by
hand in \S\ref{sec:toy} -- give \lean{toy\_certificate}: $\sigma$ is optimal against every bijection of
$\mathrm{Fin}\,6$, and \lean{toy\_cost\_eq}: the cost is exactly $7/4$. The broken potential
$\psi_{c'}=\tfrac12$ is proved infeasible (\lean{broken\_infeasible}); two further variants -- an
incorrect claimed cost, and a strict instead of non-strict inequality in the certificate's conclusion --
fail to compile. Eight theorems, no \lean{sorry}, standard axiom footprint
$\{\texttt{propext},\allowbreak\texttt{Classical.choice},\allowbreak\texttt{Quot.sound}\}$, accepted by
both the Lean kernel and the independent \texttt{nanoda} kernel~\cite{Comparator}.

\begin{keybox}
\textbf{Scope, stated plainly.} The seven real pairs use $60\times60$ floating-point cost matrices from
simulation data; re-verifying those specific numbers inside the Lean kernel is not attempted. What is
proved is the general theorem that makes any such certificate valid in principle, checked end-to-end on
the one instance small and exact enough to write down by hand and decide in the kernel.
\end{keybox}

\section{Discussion}

Fourteen of fourteen applicable pre-registered lines pass (Tables~\ref{tab:p1}--\ref{tab:certificates},
P6). That is not the headline. The headline is that \emph{passing} was never in doubt for P1 -- it is a
theorem -- and the actual content of this round is in what P2 and P4 add: a stability guarantee can be
true on every one of seven real pairs and still be the wrong tool for six of them, because its margin
scales with $\varepsilon$, not with the size of the features it is meant to protect. That is measurable,
and now measured, rather than argued. It is also consistent with, and gives a mechanism for, this
project's D1 result~\cite{Callens2026TDA}: a threshold-based detector inherits whatever margin the
stability bound offers at the resolution in play, and at $\xi/\Delta x=4$ vs.\ $8$ that margin covers
the entire diagram.

The duality certificate is the smaller but more exportable result: any pipeline computing a Wasserstein
distance between finite point sets can be made to prove its own answer, at the cost of one extra LP
solve, rather than trusting one solver call. The general theorem behind it is elementary and proved in
full; what it does not do is certify the specific $60\times60$ instances used here, which is a
statement about effort spent, not about what is possible.

The S1 correction is offered as data about method, in the same register as this project's other
retractions and null results~\cite{Callens2026TDA,Callens2026Villani}: a pre-registered pipeline with a
hand-verified toy control and negative controls still transmitted a corrupted array between two stages
without any single check catching it. The check that did catch it was neither pre-registered nor
automated -- an independent, from-scratch recomputation of one pair's numbers. That is worth stating
because it is a limit of pre-registration as practiced in this project so far, not a limit of the
mathematics: no fixed set of controls, agreed in advance, is guaranteed to catch every way a pipeline can
be wrong, and the discipline that closes that gap is closer to ``redo it independently'' than to ``check
it against a list.''

\paragraph{Availability.} Source, the pre-registration with its dated amendment, the full results
document, and every script: \url{https://github.com/xaviercallens/SocrateAI-Scientific-QuantumFluids}
(\texttt{docs/designs/CLOSED\_LOOP\_PREREG.md}, \texttt{CLOSED\_LOOP\_RESULTS.md},
\texttt{exploration/tda/loop\_*.py}, \lean{lean\_src/WassersteinCertificate.lean}).

\bibliographystyle{unsrt}
\bibliography{refs_closed_loop}
\end{document}
